\section{Control Theory Fundamentals}

\begin{purposebox}
\textbf{Purpose:} Demonstrate understanding of discrete-time dynamics, affine nonlinear systems, tracking control, and optimal control via the Hamilton--Jacobi--Bellman equation.

\textbf{When complete:} You can explain the full control pipeline from system model to optimal control law, and why each design choice was made.
\end{purposebox}

%──────────────────────────────────────────────────
\subsection{Discrete-Time Dynamic Systems}
%──────────────────────────────────────────────────

\subsubsection*{Concept Check}

\question{State-Space Representation}
Write the paper's system model
\[
  x(k+1) = x(k) + \tau\bigl[f(x(k)) + g(x(k))\,u(k)\bigr]
\]
and identify each term:
\begin{itemize}[nosep]
  \item What is $x(k) \in \mathbb{R}^n$ and what does it represent physically?
  \item What is $u(k) \in \mathbb{R}^m$ and what does it represent?
  \item What is the difference between $f(x)$ and $g(x)$?
  \item What is $\tau$?
\end{itemize}

\begin{answerbox}\end{answerbox}

\question{Time Step Selection}
The paper requires $0 < \tau < 2|x(k)|/h(x(k))$ for stability. Explain in your own words why $\tau$ cannot be arbitrary. What happens physically if $\tau$ is too large? What if $\tau$ is too small?

\begin{answerbox}\end{answerbox}

\question{Worked Example}
Consider a scalar system ($n=1$) with $f(x) = -x$, $g(x) = 1$, $\tau = 0.1$. Starting from $x(0) = 5$ with constant input $u(k) = 0$, compute $x(1)$, $x(2)$, $x(3)$. Is the state converging?

\begin{largeanswer}\end{largeanswer}

%──────────────────────────────────────────────────
\subsection{Affine Nonlinear Systems}
%──────────────────────────────────────────────────

\subsubsection*{Concept Check}

\question{Affine-in-the-Input}
The continuous-time form is $\dot{x} = f(x) + g(x)\,u$. Explain precisely what ``affine in the input'' means. Why is the control input $u$ appearing \emph{linearly} important for solving the control problem, even though $f(x)$ and $g(x)$ can be arbitrarily nonlinear?

\begin{answerbox}\end{answerbox}

\question{Affine vs.\ Linear}
Distinguish between:
\begin{itemize}[nosep]
  \item A \textbf{linear} system: $\dot{x} = Ax + Bu$
  \item An \textbf{affine} system: $\dot{x} = f(x) + g(x)\,u$
  \item A general \textbf{nonlinear} system: $\dot{x} = F(x, u)$
\end{itemize}
Why can't the paper's methods apply to the general nonlinear case?

\begin{answerbox}\end{answerbox}

\question{Controllability}
The paper requires $g(x) \neq 0$. Explain what would happen if $g(x) = 0$ for some state $x$. Why would the system be ``uncontrollable'' at that point?

\begin{answerbox}\end{answerbox}

%──────────────────────────────────────────────────
\subsection{Tracking Control Problem}
%──────────────────────────────────────────────────

\subsubsection*{Concept Check}

\question{Tracking Error}
Define the tracking error $e(k) = x(k) - r(k)$. In the ACC (Adaptive Cruise Control) context, what does $x(k)$ represent and what does $r(k)$ represent? What does $e(k) = 0$ mean physically?

\begin{answerbox}\end{answerbox}

\question{Control Law Derivation}
The paper derives the control law:
\[
  u(k) = g^{-1}(x(k))\bigl[-\beta\,\operatorname{sgn}(x(k) - r(k)) + \Delta r(k) - f(x(k))\bigr]
\]
Break this equation apart term by term:
\begin{itemize}[nosep]
  \item What does $g^{-1}(x(k))$ do?
  \item What does $-\beta\,\operatorname{sgn}(x(k) - r(k))$ do?
  \item What does $\Delta r(k)$ contribute?
  \item What does $-f(x(k))$ accomplish?
\end{itemize}

\begin{answerbox}\end{answerbox}

\question{Closed-Loop Tracking Equation}
Substituting the control law into the system model yields the tracking equation:
\[
  x(k+1) = x(k) - \beta\,\operatorname{sgn}(x(k) - r(k)) + \Delta r(k)
\]
Verify this by substituting $u(k)$ into $x(k+1) = x(k) + \tau[f(x(k)) + g(x(k))\,u(k)]$ step by step.

\begin{largeanswer}\end{largeanswer}

%──────────────────────────────────────────────────
\subsection{Optimal Control \& HJB Equation}
%──────────────────────────────────────────────────

\subsubsection*{Concept Check}

\question{Value Function}
Define the value function $V(x)$ in optimal control. What does it represent? Why is it useful?

\begin{answerbox}\end{answerbox}

\question{HJB Equation}
Write out the Hamilton--Jacobi--Bellman equation:
\[
  Q(x) + V_x^\top f(x) - \tfrac{1}{4}\,V_x^\top g(x)\,R^{-1}\,g^\top(x)\,V_x = 0
\]
Identify the role of each term:
\begin{itemize}[nosep]
  \item $Q(x)$: ?
  \item $V_x^\top f(x)$: ?
  \item The quadratic term: ?
\end{itemize}

\begin{answerbox}\end{answerbox}

\question{Analytical vs.\ Numerical}
The paper claims to derive an \emph{analytical} solution to the HJB equation. Why is this significant? What do most other approaches do instead, and what are the downsides?

\begin{answerbox}\end{answerbox}

%──────────────────────────────────────────────────
\subsection*{Self-Assessment Checklist}
%──────────────────────────────────────────────────

\begin{itemize}[leftmargin=2em, label=$\square$]
  \item I can write and interpret the discrete-time system model equation
  \item I can explain why the system is ``affine in the input'' and why that matters
  \item I can define tracking error and explain the tracking control objective
  \item I can break down the control law $u(k)$ term by term
  \item I can derive the closed-loop tracking equation by substitution
  \item I can state the HJB equation and explain each term's role
  \item I can explain why an analytical HJB solution is a significant contribution
\end{itemize}

\begin{notesbox}\end{notesbox}
